% p179_1; source PDF page 179; manually verified editable table.
\begin{tikzpicture}[x=1cm,y=-1cm,inner sep=0pt,outer sep=0pt]
\path[use as bounding box] (0,0) rectangle (16.000,8.300);
\draw[line width=.4pt] (0.000,0.000) rectangle (2.400,0.800);
\node[align=center,text width=2.140cm,font={\linespread{1}\fontsize{10}{13.50}\selectfont\bfseries}] at (1.200,0.400) {基本原理};
\draw[line width=.4pt] (2.400,0.000) rectangle (6.500,0.800);
\node[align=center,text width=3.840cm,font={\linespread{1}\fontsize{10}{13.50}\selectfont\bfseries}] at (4.450,0.400) {处理办法};
\draw[line width=.4pt] (6.500,0.000) rectangle (16.000,0.800);
\node[align=center,text width=9.240cm,font={\linespread{1}\fontsize{10}{13.50}\selectfont\bfseries}] at (11.250,0.400) {范例};
\draw[line width=.4pt] (0.000,0.800) rectangle (2.400,8.300);
\node[align=left,text width=2.140cm,font={\linespread{1}\fontsize{10}{13.50}\selectfont}] at (1.200,4.550) {为了保障人权，禁止重复处罚（评价）};
\draw[line width=.4pt] (2.400,0.800) rectangle (6.500,4.100);
\node[align=left,text width=3.840cm,font={\linespread{1}\fontsize{10}{13.50}\selectfont}] at (4.450,2.450) {禁止重复处罚行为。一个行为，触犯两个罪，只能定一个罪，不能数罪并罚};
\draw[line width=.4pt] (6.500,0.800) rectangle (16.000,4.100);
\node[align=left,text width=9.240cm,font={\linespread{1}\fontsize{10}{13.50}\selectfont}] at (11.250,2.450) {想象竞合犯、法条竞合、结果加重犯。为何对想象竞合犯要择一重罪论处？因为想象竞合犯是指一个行为同时触犯两个罪名，若数罪并罚，则意味着对一个行为进行了重复处罚，侵犯了人权};
\draw[line width=.4pt] (2.400,4.100) rectangle (6.500,8.300);
\node[align=left,text width=3.840cm,font={\linespread{1}\fontsize{10}{13.50}\selectfont}] at (4.450,6.200) {禁止重复处罚结果。两个行为，制造一个结果，对这一个结果只能处罚一次};
\draw[line width=.4pt] (6.500,4.100) rectangle (16.000,8.300);
\node[align=left,text width=9.240cm,font={\linespread{1}\fontsize{10}{13.50}\selectfont}] at (11.250,6.200) {甲杀乙，致乙重伤昏迷，以为乙死亡，为了毁尸灭迹，将乙扔进河里，乙溺亡。杀害行为和扔河里共同导致死亡，二因一果，前行为构成故意杀人罪既遂，后行为构成过失致人死亡罪。但对此不能数罪并罚，否则对一个死亡结果处罚了两次。因此，只能重罪吸收轻罪，最终定故意杀人罪既遂（多数说）};
\end{tikzpicture}
